\section*{Entry 001: Viral HRSM branch initiated}

\textbf{Date:} \today

\textbf{Purpose:} To test whether viral persistence, rebound, immune escape, and lineage survival can be represented as movement through a four-axis adaptive state space.

\[
\phi_v(t) = [H_v(t), R_v(t), S_v(t), M_v(t)].
\]

\textbf{Phase order:}
Phase 0, dataset audit.
Phase I, SARS-CoV-2 persistent infection.
Phase II, HIV longitudinal memory.
Phase III, HCV recurrence.
Phase IV, influenza antigenic drift.
Phase V, phage experimental evolution.

\textbf{Immediate decision:} Begin with SARS-CoV-2 persistent infection only.

\textbf{Failure guard:} Do not allow \(M_v\) to collapse into ordinary sequence divergence, sampling duration, or treatment labels.

\section*{Entry 002: Phase I SARS-CoV-2 candidate registry and accession audit}

\textbf{Date:} \today

\textbf{Purpose:} To replace placeholder viral candidates with a ranked SARS-CoV-2 persistent-infection registry and to separate biological relevance from actual data usability.

\textbf{Outputs generated:}
\begin{itemize}
    \item \texttt{metadata/dataset\_candidate\_registry.csv}
    \item \texttt{metadata/dataset\_candidate\_scores.csv}
    \item \texttt{metadata/sarscov2\_accession\_audit.csv}
    \item \texttt{metadata/phase1\_sarscov2\_ranked\_candidates.csv}
    \item \texttt{docs/phase1\_sarscov2\_workplan.md}
    \item \texttt{docs/phase1\_sarscov2\_workplan.tex}
\end{itemize}

\textbf{Current ranking:}
\begin{enumerate}
    \item Kemp et al., chronic infection under convalescent plasma, primary Phase I test case.
    \item Schmidt et al., 521-day persistent infection, long-memory validation case.
    \item Weigang et al., immunosuppressed within-host evolution, immune-escape validation case.
    \item Harari et al., chronic infection adaptive evolution drivers, cross-case context and robustness.
    \item Ghafari et al., community surveillance persistent infections, population background reference.
\end{enumerate}

\textbf{Interpretation:} The viral HRSM branch has moved from speculative framing to a staged dataset-audit structure. Kemp et al. is the leading Phase I opener because it offers dense longitudinal sequencing and treatment pressure. Schmidt et al. is the strongest long-memory case because the infection duration is extreme, although the reported timepoint count is smaller.

\textbf{Failure guard:} No model should be run until accession availability and sample-level structure are confirmed. Dataset priority is not equivalent to usable data.


\section*{Entry 003: First real data-contact script for Schmidt et al. ENA accession}

\textbf{Date:} \today

\textbf{Purpose:} To test real data contact using the Schmidt et al. long-memory SARS-CoV-2 persistent infection accession \texttt{PRJEB77414}.

\textbf{Script added:}
\begin{itemize}
    \item \texttt{scripts/phase1\_sarscov2/11\_fetch\_schmidt\_ena\_metadata.py}
\end{itemize}

\textbf{Expected outputs:}
\begin{itemize}
    \item \texttt{metadata/schmidt\_prjeb77414\_ena\_read\_run.tsv}
    \item \texttt{metadata/schmidt\_prjeb77414\_ena\_audit\_summary.csv}
\end{itemize}

\textbf{Interpretation rule:} This step retrieves metadata and file links only. It does not download FASTQ files. The aim is to verify whether the accession exposes run-level sequence data suitable for Phase I HRSM reconstruction.

\textbf{Failure guard:} If the accession returns no run-level files, Schmidt remains a conceptual long-memory candidate but cannot be used as the first computational case unless supplementary data can be recovered elsewhere.


\section*{Entry 003: First real data-contact script for Schmidt et al. ENA accession}

\textbf{Date:} \today

\textbf{Purpose:} To test real data contact using the Schmidt et al. long-memory SARS-CoV-2 persistent infection accession \texttt{PRJEB77414}.

\textbf{Script added:}
\begin{itemize}
    \item \texttt{scripts/phase1\_sarscov2/11\_fetch\_schmidt\_ena\_metadata.py}
\end{itemize}

\textbf{Expected outputs:}
\begin{itemize}
    \item \texttt{metadata/schmidt\_prjeb77414\_ena\_read\_run.tsv}
    \item \texttt{metadata/schmidt\_prjeb77414\_ena\_audit\_summary.csv}
\end{itemize}

\textbf{Interpretation rule:} This step retrieves metadata and file links only. It does not download FASTQ files. The aim is to verify whether the accession exposes run-level sequence data suitable for Phase I HRSM reconstruction.

\textbf{Failure guard:} If the accession returns no run-level files, Schmidt remains a conceptual long-memory candidate but cannot be used as the first computational case unless supplementary data can be recovered elsewhere.


\section*{Entry 004: Schmidt ENA accession returned no run-level files}

\textbf{Date:} \today

\textbf{Purpose:} To record the first real data-contact result for the Schmidt et al. 521-day persistent SARS-CoV-2 infection candidate.

\textbf{Result:} The ENA Portal \texttt{read\_run} query for \texttt{PRJEB77414} returned zero run-level rows, zero unique runs, zero unique samples, and no FASTQ, submitted, or BAM file links.

\textbf{Files generated:}
\begin{itemize}
    \item \texttt{metadata/schmidt\_prjeb77414\_ena\_read\_run.tsv}
    \item \texttt{metadata/schmidt\_prjeb77414\_ena\_audit\_summary.csv}
    \item \texttt{metadata/sarscov2\_accession\_audit\_updated.csv}
\end{itemize}

\textbf{Interpretation:} Schmidt et al. remains a scientifically valuable long-memory SARS-CoV-2 case because the infection duration is extreme and the paper reports five whole-genome sequencing timepoints. However, the accession does not currently expose run-level ENA files through the tested endpoint. It should therefore be retained as biological context until a usable supplementary or repository-level data path is recovered.

\textbf{Decision:} Pivot the first computational Phase I data-contact attempt back to Kemp et al., which remains the strongest dense longitudinal candidate.

\textbf{Failure guard:} This negative result is retained. It prevents accidental overclaiming that Schmidt is immediately downloadable when the tested ENA endpoint does not support that claim.


\section*{Entry 005: Kemp et al. Phase I data availability audit scaffold}

\textbf{Date:} \today

\textbf{Purpose:} To pivot from the Schmidt ENA negative result to the Kemp et al. dense longitudinal SARS-CoV-2 chronic infection candidate.

\textbf{Script added:}
\begin{itemize}
    \item \texttt{scripts/phase1\_sarscov2/13\_audit\_kemp\_data\_availability.py}
\end{itemize}

\textbf{Files generated:}
\begin{itemize}
    \item \texttt{metadata/kemp\_data\_availability\_checklist.csv}
    \item \texttt{metadata/kemp\_data\_availability\_summary.md}
    \item \texttt{results/phase1\_sarscov2/kemp\_case\_table\_template.csv}
\end{itemize}

\textbf{Known anchors:} Kemp et al. reports whole-genome ultra-deep sequencing across 23 timepoints spanning 101 days in an immunosuppressed SARS-CoV-2 infection under remdesivir and convalescent plasma pressure. The major dynamic escape genotype includes Spike D796H and Spike \(\Delta H69/\Delta V70\).

\textbf{HRSM relevance:} This case is suitable for first-pass viral HRSM because \(H_v\) can track mutational breadth, \(S_v\) can track coherence of the escape genotype, \(M_v\) can track disappearing and re-emerging genotype structure, and \(R_v\) may be recoverable from Ct, viral load, culture status, or treatment-response metadata if available.

\textbf{Failure guard:} Kemp remains the biological Phase I opener, but computational modelling must wait until sample-day mapping and sequence or variant-frequency data are recovered.


\section*{Entry 006: Kemp et al. concrete data anchors recorded}

\textbf{Date:} \today

\textbf{Purpose:} To move Kemp et al. from a biological Phase I candidate to a concrete data-anchor candidate.

\textbf{Script added:}
\begin{itemize}
    \item \texttt{scripts/phase1\_sarscov2/14\_record\_kemp\_data\_anchors.py}
\end{itemize}

\textbf{Files generated:}
\begin{itemize}
    \item \texttt{metadata/kemp\_data\_anchors.csv}
    \item \texttt{metadata/kemp\_expected\_sra\_samples.csv}
    \item \texttt{metadata/kemp\_data\_anchor\_summary.md}
\end{itemize}

\textbf{Recorded anchors:} The Kemp et al. data availability trail identifies NCBI BioProject \texttt{PRJNA682013}, expected SRA sample accessions \texttt{SAMN16976824} through \texttt{SAMN16976846}, and the GitHub repository \texttt{Steven-Kemp/sequence\_files} for short reads and figure-construction data.

\textbf{Interpretation:} Kemp et al. is now upgraded from a biological candidate to a concrete data-anchor candidate. The next step remains metadata-only: query run-level accessions, file names, file sizes, and repository contents before downloading any sequence files.

\textbf{Failure guard:} Do not download large FASTQ or long-read files until run-level metadata and sample-day mapping are clear.


\section*{Entry 007: Kemp GitHub repository inventory}

\textbf{Date:} \today

\textbf{Purpose:} To inventory the \texttt{Steven-Kemp/sequence\_files} repository before downloading any files.

\textbf{Script added:}
\begin{itemize}
    \item \texttt{scripts/phase1\_sarscov2/15\_inventory\_kemp\_github.py}
\end{itemize}

\textbf{Expected files generated:}
\begin{itemize}
    \item \texttt{metadata/kemp\_github\_sequence\_files\_inventory.csv}
    \item \texttt{metadata/kemp\_github\_hrsm\_candidate\_files.csv}
    \item \texttt{metadata/kemp\_github\_inventory\_summary.md}
\end{itemize}

\textbf{Interpretation rule:} This step is metadata-only. It identifies candidate table, variant-frequency, sample-time, abundance-proxy, and figure-source files. It does not download sequence reads or large compressed files.

\textbf{Failure guard:} Candidate file names are not evidence. A file becomes usable only after its contents are inspected and mapped to sample days, viral mutations, or HRSM axis proxies.


\section*{Entry 007: Kemp GitHub repository inventory}

\textbf{Date:} \today

\textbf{Purpose:} To inventory the \texttt{Steven-Kemp/sequence\_files} repository before downloading any files.

\textbf{Script added:}
\begin{itemize}
    \item \texttt{scripts/phase1\_sarscov2/15\_inventory\_kemp\_github.py}
\end{itemize}

\textbf{Files generated:}
\begin{itemize}
    \item \texttt{metadata/kemp\_github\_sequence\_files\_inventory.csv}
    \item \texttt{metadata/kemp\_github\_hrsm\_candidate\_files.csv}
    \item \texttt{metadata/kemp\_github\_inventory\_summary.md}
\end{itemize}

\textbf{Result:} The repository inventory found 54 files and 3 HRSM candidate files. The strongest immediate target is \texttt{figure\_data/Rawdata\_forfigures.xlsx}, a small table-like file. The compressed FASTA file under \texttt{vaccinepaper/} is deferred because it is not clearly the chronic-infection case table.

\textbf{Decision:} Download and inspect only small table-like files next. Do not download compressed read, alignment, or background files yet.

\textbf{Failure guard:} A GitHub file name is only a lead. A file becomes usable only after its contents are inspected and mapped to sample days, mutations, Ct values, treatment timing, or HRSM proxy variables.


\section*{Entry 008: Kemp small candidate files downloaded and inspected}

\textbf{Date:} \today

\textbf{Purpose:} To download only small table-like files from the Kemp GitHub inventory and inspect workbook structure without touching compressed sequence or read files.

\textbf{Script added:}
\begin{itemize}
    \item \texttt{scripts/phase1\_sarscov2/16\_download\_kemp\_small\_candidate\_tables.py}
\end{itemize}

\textbf{Expected files generated:}
\begin{itemize}
    \item \texttt{metadata/kemp\_downloaded\_small\_files.csv}
    \item \texttt{metadata/kemp\_xlsx\_sheet\_inventory.csv}
    \item \texttt{data/interim/kemp\_github\_small/}
    \item \texttt{data/interim/kemp\_github\_small/xlsx\_previews/}
\end{itemize}

\textbf{Interpretation rule:} The workbook previews determine whether Kemp can move directly to a time-indexed HRSM proxy table. If sample-day mapping, mutation frequencies, Ct values, or treatment phases appear in the workbook, full FASTQ processing can remain deferred.

\textbf{Failure guard:} Do not infer HRSM axes from figure-source tables until sheet names, columns, and row contents are inspected.


\section*{Entry 009: Kemp Fig. 5a figure-source proxy table extracted}

\textbf{Date:} \today

\textbf{Purpose:} To extract the first usable time-indexed SARS-CoV-2 HRSM proxy table from the Kemp et al. figure-source workbook.

\textbf{Script added:}
\begin{itemize}
    \item \texttt{scripts/phase1\_sarscov2/17\_extract\_kemp\_fig5a\_proxy\_features.py}
\end{itemize}

\textbf{Files generated:}
\begin{itemize}
    \item \texttt{data/interim/kemp\_fig5a\_raw.csv}
    \item \texttt{data/processed/kemp\_fig5a\_proxy\_features.csv}
    \item \texttt{results/phase1\_sarscov2/kemp\_fig5a\_hrsm\_state\_table.csv}
    \item \texttt{metadata/kemp\_fig5a\_proxy\_audit.csv}
\end{itemize}

\textbf{Interpretation:} The \texttt{Fig 5a} sheet contains 23 rows with mutation-frequency-like columns for \texttt{D796H}, \texttt{del69/70}, \texttt{W64G}, \texttt{P330S}, \texttt{Y200H}, \texttt{T240I}, and \texttt{CT}. This supports a first figure-source reconstruction of \(H_v\), \(R_v\), \(S_v\), and \(M_v\).

\textbf{Scientific caution:} This is not yet a raw sequencing analysis. It is a figure-source HRSM pilot table. The result can guide axis design, but any manuscript must label it as figure-source reconstruction unless raw variant calls or reads are processed later.

\textbf{Failure guard:} Do not run predictive modelling from this single persistent case. The output is for trajectory geometry and axis debugging only.


\section*{Entry 010: Kemp Fig. 5a proxy table cleaned}

\textbf{Date:} \today

\textbf{Purpose:} To correct the initial Fig. 5a extraction by removing blank/footer rows and fixing the missing \texttt{case\_id} assignment.

\textbf{Correction:} The first extraction retained 23 rows, but the final three rows were not biological timepoints. They were blank/footer rows from the figure-source workbook. The revised extraction retains only rows with mutation data and removes the non-biological footer rows.

\textbf{Interpretation:} The cleaned table is now suitable for trajectory plotting and axis-debugging. It remains a figure-source reconstruction, not a raw sequencing analysis.

\textbf{Failure guard:} Any HRSM values computed before this correction should be treated as deprecated because footer rows contaminated the standardization.


\section*{Entry 011: Kemp Fig. 5a HRSM trajectory plotted}

\textbf{Date:} \today

\textbf{Purpose:} To visualize the cleaned Kemp Fig. 5a figure-source HRSM trajectory across ordered mutation and Ct timepoints.

\textbf{Script added:}
\begin{itemize}
    \item \texttt{scripts/phase1\_sarscov2/18\_plot\_kemp\_fig5a\_hrsm\_trajectory.py}
\end{itemize}

\textbf{Figures generated:}
\begin{itemize}
    \item \texttt{figures/phase1\_sarscov2/kemp\_fig5a\_hrsm\_axes\_over\_time.png}
    \item \texttt{figures/phase1\_sarscov2/kemp\_fig5a\_escape\_mutations\_over\_time.png}
    \item \texttt{figures/phase1\_sarscov2/kemp\_fig5a\_phi\_over\_time.png}
    \item \texttt{figures/phase1\_sarscov2/kemp\_fig5a\_memory\_vs\_stability.png}
    \item \texttt{figures/phase1\_sarscov2/kemp\_fig5a\_reserve\_vs\_memory.png}
\end{itemize}

\textbf{Interpretation rule:} These plots are trajectory diagnostics from a figure-source reconstruction. They can be used to debug viral HRSM axis design and visualize the escape-memory structure, but they are not yet a final raw sequencing result.

\textbf{Failure guard:} No predictive claims should be made from this single-case pilot. The value here is geometric inspection and non-Markovian axis design.


\section*{Entry 012: Kemp Fig. 5a transition states classified}

\textbf{Date:} \today

\textbf{Purpose:} To classify the cleaned Kemp Fig. 5a HRSM trajectory into provisional viral transition states.

\textbf{Script added:}
\begin{itemize}
    \item \texttt{scripts/phase1\_sarscov2/19\_classify\_kemp\_fig5a\_transition\_states.py}
\end{itemize}

\textbf{Files generated:}
\begin{itemize}
    \item \texttt{results/phase1\_sarscov2/kemp\_fig5a\_transition\_states.csv}
    \item \texttt{metadata/kemp\_fig5a\_transition\_state\_summary.csv}
\end{itemize}

\textbf{Interpretation:} The provisional state labels are designed only for trajectory inspection. The key biological target is whether the D796H plus \(\Delta H69/\Delta V70\) escape structure behaves as a retained memory state, disappearing after one interval and later returning with high \(S_v\), \(M_v\), and \(\Phi_v\).

\textbf{Failure guard:} These labels are heuristic. They are not clinical classes, not viral phenotypes, and not final model outputs. They exist to guide the next axis refinement and plotting step.


\section*{Entry 013: Chronological correction for Kemp Fig. 5a HRSM pilot}

\textbf{Date:} \today

\textbf{Purpose:} To correct the figure-source HRSM pilot after reviewer audit identified row-index memory, missing \texttt{*ETA} day anchors, CT dropout plateaus in \(R_v\), and index-based plotting.

\textbf{Corrections applied:}
\begin{itemize}
    \item Manual day anchors were assigned to \texttt{93ETA*}, \texttt{99ETA*}, \texttt{100ETA*}, and \texttt{101ETA*}.
    \item CT values were linearly interpolated over explicit day anchors before \(R_v\) proxy construction.
    \item Signed numerical zero artefacts were collapsed to exact zero.
    \item The memory proxy was changed from row-adjacent similarity to a chronological exponential-decay kernel.
    \item Trajectory plots were switched from ordered row index to days since first positive RT-PCR.
\end{itemize}

\textbf{Interpretation:} This correction removes the main structural artefacts in the pilot vector space. Any previous HRSM coordinates or plots generated before this entry are deprecated.

\textbf{Failure guard:} This remains a figure-source reconstruction. The correction improves internal mathematical consistency, but it does not convert the pilot into a raw sequencing analysis.


\section*{Entry 014: Kemp transition-state labels refined}

\textbf{Date:} \today

\textbf{Purpose:} To correct the transition-state classifier so that the first observed high-escape event is not mislabeled as a return.

\textbf{Correction:} Day 82 is now treated as a coherent escape surge because it is the first observed high D796H plus \(\Delta H69/\Delta V70\) state in the figure-source trajectory. The later reappearance beginning around day 98 is retained as the candidate memory-return event.

\textbf{Interpretation:} This improves the biological meaning of the state labels. The memory claim should be attached to re-emergence after the alternate diversity interval, not to the first observed escape surge.

\textbf{Failure guard:} These labels remain heuristic trajectory annotations and should not be described as clinical classes or validated viral phenotypes.


\section*{Entry 016: Velocity sensitivity split into all-step and calendar-day summaries}

\textbf{Date:} \today

\textbf{Purpose:} To separate true calendar-day viral trajectory velocity from same-day sequencing contrast velocity.

\textbf{Script added:}
\begin{itemize}
    \item \texttt{scripts/phase1\_sarscov2/23\_summarize\_velocity\_sensitivity.py}
\end{itemize}

\textbf{Files generated:}
\begin{itemize}
    \item \texttt{metadata/kemp\_fig5a\_velocity\_summary\_all\_steps.csv}
    \item \texttt{metadata/kemp\_fig5a\_velocity\_summary\_calendar\_only.csv}
    \item \texttt{metadata/kemp\_fig5a\_velocity\_same\_day\_flags.csv}
\end{itemize}

\textbf{Interpretation:} Same-day samples are retained as useful sequencing contrasts, but they should not be mixed with calendar-day evolutionary velocity. The all-step velocity summary captures both calendar and sub-day structure, while the calendar-only summary is the safer biological velocity diagnostic.

\textbf{Failure guard:} Manuscript language should avoid interpreting same-day \(\Delta t_{\mathrm{eff}}\) velocities as literal viral evolutionary speed. They are within-day contrast magnitudes, not calendar-time rates.


\section*{Entry 017: Velocity engine normalized and state-transition matrix added}

\textbf{Date:} \today

\textbf{Purpose:} To correct the sub-day velocity explosion and add an explicit state-transition matrix for the Kemp Fig. 5a pilot trajectory.

\textbf{Scripts updated or added:}
\begin{itemize}
    \item \texttt{scripts/phase1\_sarscov2/21\_compute\_phase\_velocity.py}
    \item \texttt{scripts/phase1\_sarscov2/24\_plot\_velocity\_profiles.py}
    \item \texttt{scripts/phase1\_sarscov2/25\_build\_state\_transition\_matrix.py}
\end{itemize}

\textbf{Correction:} Same-day sequencing contrasts are no longer divided by a synthetic \(\Delta t_{\mathrm{eff}} = 0.01\). The corrected velocity engine computes step-normalized velocities for all ordered transitions and calendar-day velocities only where true positive day intervals exist.

\textbf{Interpretation:} This preserves within-day state changes without allowing sampling density to create artificial velocity explosions. The state-transition matrix records the observed directional movement among low escape baseline, coherent escape surge, alternate diversity state, retained escape state, and high-memory escape return.

\textbf{Failure guard:} The transition probabilities are empirical frequencies from one figure-source trajectory. They are not general viral transition probabilities and should not be interpreted as population-level probabilities.


\section*{Entry 019: Stationary distribution solver corrected and Phase II mapping deferred}

\textbf{Date:} \today

\textbf{Purpose:} To correct the stationary-distribution calculation after the daily-collapsed transition matrix exposed a deterministic terminal two-state cycle.

\textbf{Correction:} The stationary distribution is now computed by solving \(\pi P = \pi\) with \(\sum_i \pi_i = 1\), rather than by power iteration. This avoids numerical distortion in periodic chains, especially the deterministic retained-escape/high-memory two-state loop.

\textbf{Interpretation:} The daily-collapsed Kemp Fig. 5a pilot identifies a terminal closed class linking \texttt{retained\_escape\_state} and \texttt{high\_memory\_escape\_return}. This should be described as an observed two-state terminal cycle in a figure-source reconstruction, not as proof of a universal or permanent viral attractor.

\textbf{Phase II note:} The same daily-collapse, transition-matrix, velocity, and memory-kernel framework can later be mapped onto Phase II multi-case or cross-lineage SARS-CoV-2 datasets. That extension is deferred until Phase I is fully consolidated.

\textbf{Failure guard:} The stationary distribution is derived from one reconstructed trajectory. It is a mathematical diagnostic of this case topology only and should not be interpreted as a population-level viral residence probability.


\section*{Entry 020: Phase I Kemp Fig. 5a summary report generated}

\textbf{Date:} \today

\textbf{Purpose:} To consolidate the cleaned Phase I Kemp figure-source HRSM pilot into a single markdown and LaTeX report before any Phase II expansion.

\textbf{Script added:}
\begin{itemize}
    \item \texttt{scripts/phase1\_sarscov2/27\_generate\_phase1\_summary\_report.py}
\end{itemize}

\textbf{Files generated:}
\begin{itemize}
    \item \texttt{docs/phase1\_kemp\_fig5a\_summary\_report.md}
    \item \texttt{docs/phase1\_kemp\_fig5a\_summary\_report.tex}
    \item \texttt{metadata/phase1\_kemp\_fig5a\_summary\_report\_audit.csv}
\end{itemize}

\textbf{Interpretation:} The Phase I pilot is now consolidated as a figure-source reconstruction. It documents the low-escape baseline, coherent escape surge, alternate diversity interval, retained escape state, and high-memory return cycle, with daily-collapsed transition matrices, stationary distribution, velocity diagnostics, potential-well summaries, and figure paths.

\textbf{Failure guard:} This report is not a final manuscript and not a raw sequencing analysis. It is a reproducibility and extraction bridge for later manuscript drafting.


\section*{Entry 021: Potential-well summary patched to use corrected velocity}

\textbf{Date:} \today

\textbf{Purpose:} To remove the obsolete same-day-inflated velocity field from the Kemp Fig. 5a potential-well diagnostics.

\textbf{Script updated:}
\begin{itemize}
    \item \texttt{scripts/phase1\_sarscov2/22\_plot\_potential\_wells.py}
\end{itemize}

\textbf{Correction:} The potential-well script now reads the corrected velocity table from \texttt{kemp\_fig5a\_phase\_velocity.csv}. It reports step-normalized velocity and calendar-day velocity separately, rather than the deprecated \texttt{velocity\_norm\_total} field.

\textbf{Files regenerated:}
\begin{itemize}
    \item \texttt{results/phase1\_sarscov2/kemp\_fig5a\_potential\_wells.csv}
    \item \texttt{metadata/kemp\_fig5a\_potential\_well\_summary.csv}
    \item \texttt{figures/phase1\_sarscov2/kemp\_fig5a\_potential\_well\_over\_days.png}
    \item \texttt{figures/phase1\_sarscov2/kemp\_fig5a\_potential\_vs\_escape\_pair.png}
    \item \texttt{figures/phase1\_sarscov2/kemp\_fig5a\_potential\_vs\_step\_velocity.png}
    \item \texttt{docs/phase1\_kemp\_fig5a\_summary\_report.md}
    \item \texttt{docs/phase1\_kemp\_fig5a\_summary\_report.tex}
\end{itemize}

\textbf{Interpretation:} Potential-well depth remains defined as \(U_v=-\Phi_v\). Velocity is now treated only as a corrected companion diagnostic, separated into ordered step-space motion and calendar-day motion.

\textbf{Failure guard:} Any potential-well summary generated before this patch should be treated as deprecated because it inherited the old same-day velocity inflation.


\section*{Entry 022: Phase I Kemp Fig. 5a pilot locked}

\textbf{Date:} \today

\textbf{Purpose:} To lock the completed Phase I Kemp Fig. 5a figure-source HRSM pilot after final consistency checks.

\textbf{Final audit state:}
\begin{itemize}
    \item \texttt{n\_hrsm\_rows = 20}
    \item \texttt{n\_transition\_rows = 20}
    \item \texttt{n\_daily\_rows = 16}
    \item \texttt{n\_figures\_listed = 9}
    \item \texttt{n\_missing\_figures = 0}
    \item Stationary top states: \texttt{retained\_escape\_state} and \texttt{high\_memory\_escape\_return}, each with probability \(0.5\).
\end{itemize}

\textbf{Interpretation:} The Phase I pilot identifies a reproducible figure-source topology: low escape baseline, coherent escape surge, alternate diversity interval, retained escape state, and high-memory escape return. The daily-collapsed transition matrix identifies a terminal two-state closed class between retained escape and high-memory return.

\textbf{Boundary:} This is a figure-source reconstruction from one case. It is suitable for hypothesis generation, method validation, and later manuscript extraction, but not for universal viral claims.


\section*{Entry 023: Public GitHub reproducibility repository initialized}

\textbf{Date:} \today

\textbf{Purpose:} To prepare the viral HRSM repository for public reproducibility, reviewer inspection, and later manuscript extraction.

\textbf{Repository:}
\begin{itemize}
    \item \texttt{git@github.com:Atalebe/viral\_hrsm\_state\_space.git}
\end{itemize}

\textbf{Reproducibility additions:}
\begin{itemize}
    \item Root \texttt{README.md} rewritten for referee-facing reproduction.
    \item \texttt{LICENSE} and \texttt{CITATION.cff} added.
    \item One-command Phase I runner added: \texttt{scripts/run\_phase1\_kemp\_pipeline.sh}.
    \item Reproducibility notes added under \texttt{docs/reproducibility/}.
    \item \texttt{Makefile} added for phase execution, report generation, inventory, and cleanup.
    \item Phase I locked-output check added: \texttt{scripts/check\_phase1\_locked\_outputs.py}.
    \item \texttt{.gitignore} revised to exclude large raw sequencing files while retaining referee-facing Phase I outputs.
\end{itemize}

\textbf{Final locked Phase I audit:}
\begin{itemize}
    \item \texttt{n\_hrsm\_rows = 20}
    \item \texttt{n\_transition\_rows = 20}
    \item \texttt{n\_daily\_rows = 16}
    \item \texttt{n\_figures\_listed = 9}
    \item \texttt{n\_missing\_figures = 0}
    \item Stationary terminal states: \texttt{retained\_escape\_state} and \texttt{high\_memory\_escape\_return}, each with probability \(0.5\).
\end{itemize}

\textbf{Failure guard:} The repository explicitly labels Phase I as a figure-source reconstruction, not a raw sequencing analysis. Raw sequencing-scale files are not committed. The public repository preserves scripts, metadata, results, figures, reports, and audit files needed to reproduce the locked pilot.


\section*{Entry 024: Phase II SARS-CoV-2 multi-case scaffold repaired and initialized}

\textbf{Date:} \today

\textbf{Purpose:} To repair and initialize Phase II as a cautious generalization test of the Phase I Kemp HRSM topology.

\textbf{Phase II scope:} Phase II will test whether the Phase I topology recurs across additional persistent SARS-CoV-2 cases or cohorts. The target is not to prove a universal viral law, but to determine whether the separation between alternate diversity and retained escape or high-memory return appears outside the Kemp figure-source reconstruction.

\textbf{Directories initialized:}
\begin{itemize}
    \item \texttt{scripts/phase2\_sarscov2\_multicase/}
    \item \texttt{data/raw/phase2\_sarscov2\_multicase/}
    \item \texttt{data/interim/phase2\_sarscov2\_multicase/}
    \item \texttt{data/processed/phase2\_sarscov2\_multicase/}
    \item \texttt{results/phase2\_sarscov2\_multicase/}
    \item \texttt{figures/phase2\_sarscov2\_multicase/}
    \item \texttt{metadata/phase2\_sarscov2\_multicase/}
    \item \texttt{docs/phase2\_sarscov2\_multicase/}
\end{itemize}

\textbf{Scripts added:}
\begin{itemize}
    \item \texttt{scripts/phase2\_sarscov2\_multicase/00\_initialize\_phase2.py}
    \item \texttt{scripts/phase2\_sarscov2\_multicase/01\_validate\_phase2\_registry.py}
\end{itemize}

\textbf{Failure guard:} Phase II must not use unrelated population consensus genomes as if they were within-host trajectories. Each case must have ordered timepoints and explicit day or date mapping.


\section*{Entry 025: Phase II accession audit initiated}

\textbf{Date:} \today

\textbf{Purpose:} To begin Phase II data-contact auditing with accession-level metadata retrieval rather than direct data download.

\textbf{Script added:}
\begin{itemize}
    \item \texttt{scripts/phase2\_sarscov2\_multicase/03\_audit\_phase2\_accessions.py}
\end{itemize}

\textbf{Initial target:}
\begin{itemize}
    \item \texttt{BOYLE\_2025\_PRJNA1295507}
\end{itemize}

\textbf{Interpretation rule:} This step queries run-level metadata and file-link availability only. It does not download raw sequence files. A candidate becomes computationally actionable only after run rows, sample mapping, file links, and longitudinal structure are confirmed.

\textbf{Failure guard:} BioProject existence alone is not sufficient for Phase II modelling. The accession audit must confirm usable run-level files and later sample-day mapping.


\section*{Entry 026: Boyle Phase II run manifest built}

\textbf{Date:} \today

\textbf{Purpose:} To convert the Boyle \texttt{PRJNA1295507} ENA run-level audit into a structured Phase II run manifest without downloading raw sequence files.

\textbf{Script added:}
\begin{itemize}
    \item \texttt{scripts/phase2\_sarscov2\_multicase/04\_build\_boyle\_run\_manifest.py}
\end{itemize}

\textbf{Files generated:}
\begin{itemize}
    \item \texttt{metadata/phase2\_sarscov2\_multicase/boyle\_run\_manifest.csv}
    \item \texttt{metadata/phase2\_sarscov2\_multicase/boyle\_run\_manifest\_summary.csv}
    \item \texttt{data/interim/phase2\_sarscov2\_multicase/boyle\_sample\_day\_mapping\_template.csv}
\end{itemize}

\textbf{Interpretation:} Boyle is now a real Phase II data-contact candidate because the BioProject returns run-level metadata and file links. It is not yet HRSM-ready because sample-day mapping, patient or infection grouping, and longitudinal structure still need to be recovered.

\textbf{Failure guard:} No FASTQ files should be downloaded until the sample-day mapping template is completed or a supplementary table provides the required patient-time structure.

\section*{Entry 029: Boyle PMC appendix and table extraction}

\textbf{Date:} \today

\textbf{Purpose:} To recover Boyle supplementary and article table metadata directly from the PMC page after the PMC binary supplement link returned an HTML download-preparation page instead of a ZIP payload.

\textbf{Script added:}
\begin{itemize}
    \item \texttt{scripts/phase2\_sarscov2\_multicase/07\_extract\_boyle\_pmc\_appendix\_tables.py}
\end{itemize}

\textbf{Files generated:}
\begin{itemize}
    \item \texttt{metadata/phase2\_sarscov2\_multicase/boyle\_pmc\_appendix\_text.txt}
    \item \texttt{metadata/phase2\_sarscov2\_multicase/boyle\_pmc\_appendix\_tables\_inventory.csv}
    \item \texttt{metadata/phase2\_sarscov2\_multicase/boyle\_pmc\_linked\_tables\_inventory.csv}
    \item \texttt{metadata/phase2\_sarscov2\_multicase/boyle\_pmc\_lc\_sample\_hits.csv}
\end{itemize}

\textbf{Interpretation:} This step attempts to recover sample identifiers, table metadata, and supplementary content directly from the PMC HTML. The immediate target is to find whether LC sample titles from the ENA manifest appear in article or appendix tables.

\textbf{Failure guard:} If LC sample IDs are absent from the PMC table text, the sample-day mapping must be recovered from the publisher supplement manually or through an alternate download path before Boyle can enter HRSM extraction.

\section*{Entry 030: Boyle BioSample metadata expanded}

\textbf{Date:} \today

\textbf{Purpose:} To query NCBI BioSample metadata for all Boyle \texttt{SAMN} sample accessions after PMC and supplement-download routes failed to expose LC sample-day mapping directly.

\textbf{Script added:}
\begin{itemize}
    \item \texttt{scripts/phase2\_sarscov2\_multicase/08\_expand\_boyle\_biosample\_metadata.py}
\end{itemize}

\textbf{Files generated:}
\begin{itemize}
    \item \texttt{metadata/phase2\_sarscov2\_multicase/boyle\_biosample\_metadata.csv}
    \item \texttt{metadata/phase2\_sarscov2\_multicase/boyle\_run\_manifest\_with\_biosample\_metadata.csv}
    \item \texttt{metadata/phase2\_sarscov2\_multicase/boyle\_biosample\_metadata\_summary.csv}
\end{itemize}

\textbf{Interpretation:} This step attempts to recover collection dates, isolate labels, host fields, and any patient/time grouping metadata from NCBI BioSample records. Boyle remains non-HRSM-ready until the sample-day or patient-time mapping is confirmed.

\textbf{Failure guard:} BioSample metadata may expose sample-level attributes but not necessarily longitudinal grouping. If collection dates or patient identifiers remain absent, Supplemental Table S1 must still be recovered through a manual or alternate route.


\section*{Entry 031: Boyle temporal metadata audited}

\textbf{Date:} \today

\textbf{Purpose:} To determine whether BioSample metadata is sufficient to convert Boyle from a run-level Phase II candidate into a patient-level longitudinal HRSM dataset.

\textbf{Script added:}
\begin{itemize}
    \item \texttt{scripts/phase2\_sarscov2\_multicase/09\_audit\_boyle\_temporal\_grouping.py}
\end{itemize}

\textbf{Files generated:}
\begin{itemize}
    \item \texttt{metadata/phase2\_sarscov2\_multicase/boyle\_temporal\_metadata\_timeline.csv}
    \item \texttt{metadata/phase2\_sarscov2\_multicase/boyle\_temporal\_grouping\_audit\_summary.csv}
    \item \texttt{metadata/phase2\_sarscov2\_multicase/boyle\_collection\_date\_counts.csv}
    \item \texttt{metadata/phase2\_sarscov2\_multicase/boyle\_candidate\_date\_series.csv}
\end{itemize}

\textbf{Interpretation:} BioSample metadata provides collection dates, host fields, geographic fields, and GISAID accession attributes for the Boyle samples. This is sufficient for date mapping but not sufficient for HRSM trajectory extraction unless patient or infection grouping is recovered.

\textbf{Failure guard:} Date order across unrelated samples must not be treated as within-host viral evolution. Boyle remains blocked for HRSM modelling until patient-level grouping is recovered from Supplemental Table S1 or another source.


\section*{Entry 032: Boyle registry status updated after temporal grouping audit}

\textbf{Date:} \today

\textbf{Purpose:} To update the Phase II registry after Boyle BioSample expansion and temporal grouping audit.

\textbf{Script added:}
\begin{itemize}
    \item \texttt{scripts/phase2\_sarscov2\_multicase/10\_update\_phase2\_registry\_after\_boyle\_audit.py}
\end{itemize}

\textbf{Interpretation:} Boyle has strong run-level and BioSample metadata. The BioProject returns 50 runs, and BioSample records provide collection dates and GISAID accessions for all 50 samples. However, no patient or infection grouping is exposed in the recovered metadata. Boyle is therefore date-mapped but not patient-grouped.

\textbf{Decision:} Boyle remains blocked for HRSM trajectory extraction until Supplemental Table S1 or another source provides patient/infection IDs and longitudinal grouping.

\textbf{Failure guard:} The 50 Boyle samples must not be treated as one chronological within-host trajectory. Date order across unrelated samples is not viral memory.

\section*{Entry 033: Chaguza Phase II accession audit initiated}

\textbf{Date:} \today

\textbf{Purpose:} To move to the next Phase II candidate after Boyle was blocked by missing patient/infection grouping.

\textbf{Script added:}
\begin{itemize}
    \item \texttt{scripts/phase2\_sarscov2\_multicase/11\_audit\_chaguza\_accessions.py}
\end{itemize}

\textbf{Rationale:} Chaguza et al. is a long-duration persistent SARS-CoV-2 infection candidate and is closer to the Phase I Kemp problem than a broad sample panel. The immediate goal is accession and table discovery, not modelling.

\textbf{Files generated:}
\begin{itemize}
    \item \texttt{metadata/phase2\_sarscov2\_multicase/chaguza\_article\_accession\_summary.csv}
    \item \texttt{metadata/phase2\_sarscov2\_multicase/chaguza\_article\_links.csv}
    \item \texttt{metadata/phase2\_sarscov2\_multicase/chaguza\_accession\_candidates.csv}
    \item \texttt{metadata/phase2\_sarscov2\_multicase/chaguza\_ena\_accession\_audit.csv}
\end{itemize}

\textbf{Failure guard:} Chaguza should only proceed if accession-level files and longitudinal sample mapping can be recovered. A long infection duration alone is not sufficient for HRSM extraction.


\section*{Entry 034: Weigang Phase II accession audit initiated}

\textbf{Date:} \today

\textbf{Purpose:} To audit Weigang et al. as the next Phase II candidate after Boyle was blocked by missing patient grouping and Chaguza exposed no accession candidates in the PMC article text.

\textbf{Script added:}
\begin{itemize}
    \item \texttt{scripts/phase2\_sarscov2\_multicase/12\_audit\_weigang\_accessions.py}
\end{itemize}

\textbf{Rationale:} Weigang et al. is a longitudinal immunosuppressed-patient SARS-CoV-2 within-host evolution case. It is therefore closer to the Phase I Kemp trajectory problem than broad panel datasets.

\textbf{Files generated:}
\begin{itemize}
    \item \texttt{metadata/phase2\_sarscov2\_multicase/weigang\_article\_accession\_summary.csv}
    \item \texttt{metadata/phase2\_sarscov2\_multicase/weigang\_article\_links.csv}
    \item \texttt{metadata/phase2\_sarscov2\_multicase/weigang\_accession\_candidates.csv}
    \item \texttt{metadata/phase2\_sarscov2\_multicase/weigang\_ena\_accession\_audit.csv}
\end{itemize}

\textbf{Failure guard:} Weigang should proceed only if accession-level files, longitudinal sample mapping, or extractable mutation-frequency tables are recovered.

\section*{Entry 035: Weigang run manifest and BioSample metadata expanded}

\textbf{Date:} \today

\textbf{Purpose:} To convert the Weigang \texttt{ERP132087} accession audit into a structured run manifest and BioSample metadata table.

\textbf{Scripts added:}
\begin{itemize}
    \item \texttt{scripts/phase2\_sarscov2\_multicase/13\_build\_weigang\_run\_manifest.py}
    \item \texttt{scripts/phase2\_sarscov2\_multicase/14\_expand\_weigang\_biosample\_metadata.py}
\end{itemize}

\textbf{Interpretation:} Weigang now has run-level ENA access through \texttt{ERP132087}. The immediate question is whether the 12 samples expose enough date, title, or timepoint metadata to reconstruct a single-patient longitudinal trajectory.

\textbf{Failure guard:} No raw FASTQ files should be downloaded until sample-day mapping is recovered or the BioSample/sample-title metadata clearly encode the longitudinal order.


\section*{Entry 036: Weigang longitudinal trajectory mapped}

\textbf{Date:} \today

\textbf{Purpose:} To convert Weigang from a run-level accession success into the first mapped Phase II longitudinal SARS-CoV-2 HRSM candidate.

\textbf{Scripts added:}
\begin{itemize}
    \item \texttt{scripts/phase2\_sarscov2\_multicase/15\_build\_weigang\_longitudinal\_map.py}
    \item \texttt{scripts/phase2\_sarscov2\_multicase/16\_update\_phase2\_registry\_after\_weigang\_mapping.py}
\end{itemize}

\textbf{Interpretation:} Weigang sample titles encode longitudinal days. The swab samples define the primary within-host trajectory, while isolate samples are retained as controls. This makes Weigang the first true Phase II mapping candidate after Boyle was blocked by missing patient grouping and Chaguza exposed no accession candidates.

\textbf{Failure guard:} Weigang is mapping-ready, not raw-feature-ready. HRSM extraction should wait until mutation-frequency or raw-derived feature construction is planned explicitly.

\section*{Entry 037: Weigang metadata-only HRSM scaffold built}

\textbf{Date:} \today

\textbf{Purpose:} To create a metadata-only scaffold for the mapped Weigang primary swab trajectory before any raw-derived feature extraction.

\textbf{Script added:}
\begin{itemize}
    \item \texttt{scripts/phase2\_sarscov2\_multicase/17\_build\_weigang\_metadata\_hrsm\_scaffold.py}
\end{itemize}

\textbf{Files generated:}
\begin{itemize}
    \item \texttt{data/processed/phase2\_sarscov2\_multicase/weigang\_metadata\_proxy\_features.csv}
    \item \texttt{metadata/phase2\_sarscov2\_multicase/weigang\_metadata\_hrsm\_scaffold\_audit.csv}
    \item \texttt{metadata/phase2\_sarscov2\_multicase/weigang\_raw\_feature\_readiness\_table.csv}
\end{itemize}

\textbf{Interpretation:} The Weigang trajectory is mapped across nine primary swab timepoints from day 0 to day 140. The scaffold records technical metadata and explicit placeholders for the raw-derived mutation, stability, and memory features required for real HRSM extraction.

\textbf{Failure guard:} The metadata scaffold is not a viral HRSM result. Sequencing depth must not be mistaken for biological reserve, recoverability, stability, or memory.
